\documentclass[11pt]{article}

\usepackage[T1]{fontenc}
\usepackage[utf8]{inputenc}
\usepackage{amsmath,amssymb,amsthm}
\usepackage{booktabs}
\usepackage{geometry}
\usepackage[hidelinks]{hyperref}
\usepackage{microtype}

\geometry{margin=1in}
\hypersetup{
  pdftitle={The Order-Seven Degree-Square Turan Problem for the Regular Five-Vertex Tournament},
  pdfauthor={Xiaojian Zeng},
  pdfsubject={An exact computer-assisted result in extremal digraph theory},
  pdfkeywords={degree-square Turan problem, regular tournament, exact enumeration, extremal digraph}
}

\newtheorem{theorem}{Theorem}
\newtheorem{lemma}[theorem]{Lemma}
\newcommand{\Reg}{\operatorname{Reg}_5}
\newcommand{\exout}{\operatorname{ex}^{+}_{L_2}}

\title{The Order-Seven Degree-Square Tur\'an Problem\\
for the Regular Five-Vertex Tournament}
\author{Xiaojian Zeng\\
\small PBC School of Finance, Tsinghua University, Beijing 100190, China\\
\small \texttt{xiaojianzeng@foxmail.com}}
\date{}

\begin{document}
\maketitle

\begin{abstract}
Let $\Reg$ be the regular tournament on five vertices, and allow digraphs to
contain antiparallel arcs.  An exhaustive exact computation proves that the
largest value of $\sum_v d^+(v)^2$ over $\Reg$-free digraphs on seven labeled
vertices is $186$.  There are exactly $105$ labeled maximizers, all balanced
complete bidirected four-partite graphs, with part sizes $2,2,2,1$.  The
certificate checks all $7{,}377{,}595$ digraphs whose objective is at least the
claimed value and is replayed by a differently organized enumeration of all
$6{,}937{,}603$ strict candidates.  This is an order-seven result only: the
all-order conjecture remains unresolved for $m\geq 8$.
\end{abstract}

\section{Statement and context}

A digraph here is loopless, has at most one arc in each ordered direction,
and may have both arcs between a pair of vertices.  A forbidden copy need not
be induced.  For a fixed digraph $F$, write
\[
 \exout(m,F)=\max\left\{\sum_{v\in V(D)}d_D^+(v)^2:
 |V(D)|=m,\ D\text{ is }F\text{-free}\right\}.
\]
This is the one-sided degree-square Tur\'an number used by Ai, Chen, Chen,
Yan and Zhao~\cite{AiEtAl2026}; for self-converse tournaments it also equals a
natural finite palette extremal quantity.

Identify the vertices of $\Reg$ with $\mathbb Z/5\mathbb Z$ and put an arc
$i\to j$ when $j-i\pmod 5$ belongs to $\{1,2\}$.  Thus $\Reg$ is the regular
five-vertex tournament.  Let $T_4(m)$ denote the balanced complete bidirected
four-partite graph: there are no arcs within a part and both arcs occur between
different parts.

I\v lkovi\v c~\cite[Conjecture~3.3]{Ilkovic2026} conjectured that
\[
 \exout(m,\Reg)=\lVert d^+_{T_4(m)}\rVert_2^2
 \qquad(m\geq 1).
\]
That work reports small-order computation as motivation, but does not give an
order-seven certificate or an equality-case census.  The result below supplies
such a finite certification.  It neither proves nor disproves the conjecture at
any $m\geq 8$.

\begin{theorem}\label{thm:main}
Among loopless digraphs on the labeled vertex set $\{0,\ldots,6\}$, with
antiparallel arcs allowed, every $\Reg$-free digraph $D$ satisfies
\[
  \sum_{v=0}^6 d_D^+(v)^2\leq 186.
\]
Equality holds for exactly $105$ labeled digraphs.  They are precisely the
copies of $T_4(7)$ whose part sizes are $2,2,2,1$.  In particular,
$\exout(7,\Reg)=186$.
\end{theorem}

\section{Lower bound}

The graph $T_4(7)$ has one vertex of out-degree $6$ and six vertices of
out-degree $5$.  Hence
\[
 \lVert d^+_{T_4(7)}\rVert_2^2=6^2+6\cdot5^2=186.
\]
It is $\Reg$-free: any five selected vertices contain two in the same one of
the four parts, and that pair has no arc, whereas every pair in a tournament
has an arc.  This proves the lower bound without computation.

\section{Exhaustive upper-bound certificate}

All computations described below use integer arithmetic and bit-subset tests.
The two complete programs and their saved outputs are
\texttt{reg5\_m7\_enum.py}, \texttt{reg5\_m7\_enum.log},
\texttt{replay\_reg5\_m7.py}, and \texttt{reg5\_m7\_replay.log} in the
ancillary \texttt{evidence} directory.

\subsection{Primary threshold enumeration}

The six possible outgoing arcs from each vertex form a six-bit row.  Seven
independent rows therefore give a bijection between $42$-bit words and all
$2^{42}$ loopless labeled digraphs on seven vertices.  A row of popcount $d$
contributes $d^2$ to the objective.

The depth-first traversal stores a partial score $s$ after $i$ rows.  It prunes
a branch only if
\[
 s+36(7-i)<186.\tag{1}\label{eq:prune}
\]
Since $36$ is the largest possible contribution of one remaining row, every
digraph of score at least $186$ reaches a leaf.  As an exact count check, let
\[
 P(z)=\sum_{d=0}^{6}\binom{6}{d}z^{d^2}.
\]
The coefficient of $z^t$ in $P(z)^7$ is the number of row words of score $t$.
An integer coefficient recurrence gives
\[
 [z^{186}]P(z)^7=439{,}992,
 \qquad \sum_{t>186}[z^t]P(z)^7=6{,}937{,}603.
\]
The traversal visits their sum, $7{,}377{,}595$, and its score-by-score
histogram agrees with the recurrence.

It remains to describe the forbidden-pattern test.  On each of the
$\binom75=21$ five-vertex supports, the program generates every labeling of
the cyclic representative of $\Reg$ and deduplicates the resulting arc sets.
The cyclic tournament has five automorphisms, so this produces
\[
 \binom75\frac{5!}{5}=21\cdot24=504
\]
distinct ten-arc masks.  A digraph contains $\Reg$ exactly when one of these
masks is a subset of its arc mask.  In particular, no restriction is imposed
on reverse arcs, so the test is for a not-necessarily-induced copy.

Every one of the $6{,}937{,}603$ candidates of score greater than $186$
contains a forbidden mask.  At score $186$, exactly $105$ candidates contain
no forbidden mask.  Together with the lower bound, this proves the numerical
part of Theorem~\ref{thm:main}.

\subsection{Differently organized replay}

The replay does not use the score-pruned traversal or its dense forbidden
masks.  It first enumerates the vectors
\[
 (d_0,\ldots,d_6)\in\{0,\ldots,6\}^7,
 \qquad \sum_{i=0}^6d_i^2>186.
\]
There are exactly $2{,}647$ such vectors.  For a fixed vector it takes the
Cartesian product of the $\binom6{d_i}$ row subsets having the prescribed
degrees.  Summing these product sizes gives $6{,}937{,}603$ digraphs, agreeing
with the strict coefficient count above.  Each forbidden copy is represented
instead as ten positions in a $7\times7$ adjacency array.  This replay again
finds a forbidden copy in every strict candidate.

The equality masks saved by the primary computation are then parsed rather
than regenerated.  The replay recomputes score $186$ and $\Reg$-freeness for
each of the $105$ masks.  This separation is not a formal verification of the
runtime, but it detects errors tied to the primary traversal, its dense mask
layout, or its strict-level accounting.  Similar use of exhaustive integer
search together with a separate exact replay has supported other deliberately
finite-scope results; see, for example, Zeng, Liu, Ratnavelu and Huat
\cite{ZengEtAl2026}.

\section{Equality classification}

For each of the $105$ equality masks, the replay checks every unordered pair
$\{u,v\}$.  The pair is either bidirected or has neither arc.  The graph formed
by the pairs of the second kind has exactly three edges, three components of
size two, and one isolated vertex.  Its components are therefore four
independent parts of sizes $2,2,2,1$, and all pairs in different components
are bidirected.  Thus every enumerated equality case is $T_4(7)$.

Conversely, the number of partitions of seven labeled vertices into three
unlabeled pairs and one singleton is
\[
 \frac{7!}{(2!)^3 3!}=105.
\]
Each partition defines one labeled copy of $T_4(7)$, and the lower-bound
argument shows that it is $\Reg$-free with score $186$.  Hence the equality
census is complete.

\section{Reproducibility and limitations}

From the project root, the decisive artifacts can be regenerated and checked
with
\begin{verbatim}
python3 evidence/reg5_m7_enum.py > evidence/reg5_m7_enum.log
python3 evidence/replay_reg5_m7.py > evidence/reg5_m7_replay.log
sha256sum -c evidence/reg5_m7_SHA256SUMS
\end{verbatim}
The programs use Numba to compile the exhaustive loops, but every score,
coefficient, adjacency entry, and containment comparison is integral.  There
is no random sampling, floating-point comparison, or numerical tolerance.
The trusted computational base still includes the Python and Numba runtimes,
the hardware, and faithful execution of the supplied code.

The theorem exhausts only the seven-vertex domain.  It supplies no structural
upper bound for $m\geq8$ and should not be read as evidence that finite
enumeration alone will settle the all-order conjecture.

\paragraph{Computational and AI assistance.}
OpenAI Codex assisted with the literature search, exact-search design, draft
code, adversarial audit organization, and initial manuscript preparation.
Language-model output is not used as mathematical evidence: the finite claim
rests on the explicit reduction and the released exact programs and logs.  No
proof assistant was used.

\bibliographystyle{plain}
\bibliography{references}

\end{document}
